\chapter*{Resumen}
El proyecto aborda la complejidad de diseñar experiencias d ejuego altamente rejugables y técnicamente sólidas bajo una arquitectura de Ingeniería del Software, tomando como base el género Roguelike. El problema principal reside en equilibrar la generación aleatoria de niveles con una narrativa coherente y  mecánicas de combate precisas. Para ellos, se propone el desarrollo de Via Inferni, un videojuego 2D en Unity que adapta la estructura literaria de El Infierno de Dante Alighieri. La solución implementa un sistema de generación procedural de niveles distribuidos en nueve círculos temáticos, un sistema de combate en tiempo real y una mecánica de “cambio dual” entre dos personajes (Dante y Virgilio). El resultado es una prueba de concepto modular que integra gestión de entidades, progresión por partidas (runs) y una interfaz funcional, demostrando la viabilidad de un desarrollo interactivo profesional y escalable.